% Front matter for AI Almanach.
% Keep this file independent from chapter numbering.

\begin{titlepage}
    \thispagestyle{empty}
    \begin{tikzpicture}[remember picture,overlay]
        \fill[\maincolor!8!white]
            (current page.south west) rectangle (current page.north east);
        \fill[\maincolor!75!black]
            (current page.north west) rectangle
            ([yshift=-3.2cm]current page.north east);
        \draw[\maincolor!45!white,line width=1.2pt]
            ([xshift=1.3cm,yshift=-4.3cm]current page.north west)
            -- ([xshift=-1.3cm,yshift=-4.3cm]current page.north east);
        \foreach \x/\y/\r in {
            2.0/4.0/0.18, 4.1/5.0/0.13, 6.2/3.7/0.16,
            8.0/5.4/0.12, 10.2/4.2/0.19, 12.1/5.1/0.14,
            14.0/3.8/0.17, 16.2/5.0/0.13, 18.1/4.0/0.18}
        {
            \fill[\maincolor!60!black]
                ([xshift=\x cm,yshift=\y cm]current page.south west)
                circle (\r cm);
        }
        \draw[\maincolor!35!black,line width=0.6pt]
            ([xshift=2.0cm,yshift=4.0cm]current page.south west)
            -- ([xshift=4.1cm,yshift=5.0cm]current page.south west)
            -- ([xshift=6.2cm,yshift=3.7cm]current page.south west)
            -- ([xshift=8.0cm,yshift=5.4cm]current page.south west)
            -- ([xshift=10.2cm,yshift=4.2cm]current page.south west)
            -- ([xshift=12.1cm,yshift=5.1cm]current page.south west)
            -- ([xshift=14.0cm,yshift=3.8cm]current page.south west)
            -- ([xshift=16.2cm,yshift=5.0cm]current page.south west)
            -- ([xshift=18.1cm,yshift=4.0cm]current page.south west);
    \end{tikzpicture}

    \vspace*{1.1cm}
    {\color{white}\sffamily\bfseries\fontsize{34}{38}\selectfont
        AI Almanach\par}
    \vspace{0.35cm}
    {\color{white}\sffamily\fontsize{16}{20}\selectfont
        Machine Learning, Data Science, and Intelligent Systems\par}
    \vspace{0.15cm}
    {\color{white}\sffamily\fontsize{13}{16}\selectfont
        A step-by-step engineering reference\par}

    \vfill
    \begin{tcolorbox}[
        enhanced,
        width=0.72\textwidth,
        colback=white!94!\maincolor,
        colframe=\maincolor!70!black,
        boxrule=0.6pt,
        arc=1mm]
        \sffamily
        From vectors to validated AI systems; from one gradient step to a
        tiny language model with a LoRA adapter.
        The destination is competence, not vocabulary bingo.
    \end{tcolorbox}
    \vfill

    {\sffamily\Large Bernhard Gerlach\par}
    \vspace{0.25cm}
    {\sffamily\large AI Almanach, version 2 working edition\par}
    {\sffamily\normalsize \today\par}
\end{titlepage}

\pagenumbering{roman}
\clearpage
\thispagestyle{empty}
\vspace*{\fill}
{\sffamily\small
\textbf{Purpose.}
This AI Almanach is an educational engineering reference, not a substitute for
domain-specific professional advice, safety assurance, or legal review.

\medskip
\textbf{Scientific status.}
Definitions, equations, algorithms, and empirical claims should be traceable
to cited primary literature, standards, official documentation, or clearly
identified secondary sources.
Claims about fast-moving models, software, law, and benchmarks carry a source
date and must be rechecked before consequential use.

\medskip
\textbf{Reproducibility.}
Practical results depend on data, random seeds, software versions, numerical
precision, hardware, and measurement protocol.
A benchmark without its conditions is a story, not a measurement.

\medskip
Typeset with LuaLaTeX using a card-based design built on \texttt{tcolorbox},
\texttt{TikZ}, and \texttt{biblatex}.
}
\vspace*{\fill}
\clearpage

\chapter*{Abstract}
\addcontentsline{toc}{chapter}{Abstract}

Artificial intelligence is not one technique.
It is a family of ways to represent problems, learn from data, reason under
uncertainty, optimise decisions, and build systems that interact with people
and the physical world.
The AI Almanach develops that family from first principles for engineering
students with mixed backgrounds.

The route begins with mathematical and programming refreshers, continues
through data analysis and classical machine learning, and then builds neural
networks, reinforcement learning, computer vision, natural language
processing, generative models, large language models, and multi-agent systems.
Applications are connected to the engineering constraints that determine
whether a model is useful: data quality, evaluation design, latency, memory,
energy, safety, security, privacy, maintainability, and human oversight.

The central didactic rule is simple: one concept at a time, but never in
isolation.
Every important idea is introduced through motivation, plain-language
intuition, notation and assumptions, a worked example, an engineering note,
a failure mode, a check for understanding, and a practical application.
Testing and evaluation form a continuous spine rather than a final chapter
that can be politely ignored when a deadline approaches.

\chapter*{Preface: learning AI without losing the plot}
\addcontentsline{toc}{chapter}{Preface}

AI can feel difficult for an innocent reason: several kinds of knowledge are
stacked on top of one another.
A learner may be meeting a new application, a new mathematical notation, a
new algorithm, and a new software tool in the same paragraph.
That is four lessons wearing one trench coat.

This book separates those layers and then reconnects them deliberately.
You will repeatedly move through the same engineering loop:

\begin{center}
\begin{tikzpicture}[
    node distance=8mm and 10mm,
    every node/.style={
        draw=\maincolor!65!black,
        fill=\maincolor!5!white,
        rounded corners=1mm,
        font=\sffamily\small,
        align=center,
        minimum height=8mm},
    >={Stealth[length=2mm]}]
    \node (question) {Question\\and constraints};
    \node[right=of question] (data) {Data and\\assumptions};
    \node[right=of data] (model) {Model and\\objective};
    \node[below=of model] (eval) {Evaluation\\and uncertainty};
    \node[left=of eval] (system) {System and\\operations};
    \node[left=of system] (decision) {Decision, feedback,\\and revision};
    \draw[->] (question) -- (data);
    \draw[->] (data) -- (model);
    \draw[->] (model) -- (eval);
    \draw[->] (eval) -- (system);
    \draw[->] (system) -- (decision);
    \draw[->] (decision) -- (question);
\end{tikzpicture}
\captionof{figure}{The engineering learning loop used throughout the AI Almanach.}
\label{fig:engineering-learning-loop}
\end{center}

The loop matters because model training is only one step.
An elegant model trained on a leaked target, evaluated with the wrong metric,
or deployed without monitoring is still an unsuccessful engineering system.

\chapter*{How to use the AI Almanach}
\addcontentsline{toc}{chapter}{How to use the AI Almanach}

\section*{Choose a route}

The chapters are dependency ordered, but the book is also a reference.
Choose the route that matches your immediate goal and return to prerequisites
when a symbol or assumption feels unfamiliar.

\begin{table}[htbp]
\centering
\caption{Suggested learning routes through the AI Almanach.}
\label{tab:learning-routes}
\begin{tblr}{
    width=\textwidth,
    colspec={Q[l,0.18\textwidth] X[l] X[l]},
    row{1}={font=\bfseries, bg=\maincolor!15!white},
    hlines,
    vlines}
Route & Begin with & Then emphasise \\
First complete pass & Orientation, mathematics, Python, and data &
Classical ML, deep learning, evaluation, then the domain chapters \\
ML engineer & Mathematics refresh and data foundations &
Classical ML, deep learning, evaluation/MLOps, and practical labs \\
LLM engineer & Linear algebra, probability, neural networks &
NLP, Transformers, practical LLM engineering, evaluation, safety \\
Data scientist & Statistics, Python, and data analytics &
Classical ML, causal reasoning, visualisation, and deployment \\
Technical leader & Introduction and evaluation &
Safety, law, operations, applications, and evidence limitations \\
\end{tblr}
\end{table}

\section*{The standard learning unit}

New and substantially revised material follows the same sequence.
The sequence reduces cognitive switching and makes omissions visible during
review.

\begin{enumerate}
    \item \textbf{Why it matters:} the question the concept helps answer.
    \item \textbf{Prerequisites:} the minimum earlier ideas and notation.
    \item \textbf{Plain language:} a useful mental model, including its limits.
    \item \textbf{Formal model:} definitions, shapes, units, and assumptions.
    \item \textbf{Worked example:} small enough to calculate by hand.
    \item \textbf{Engineer's note:} implementation and systems consequences.
    \item \textbf{Failure mode:} how the idea is commonly misused.
    \item \textbf{Check:} a retrieval or transfer question with a testable
          answer.
    \item \textbf{Practice:} a reproducible experiment with a baseline.
    \item \textbf{Evidence:} citations and the scope of what they support.
\end{enumerate}

\begin{engineernote}
You do not understand an algorithm merely because its name feels familiar.
You understand it when you can state its inputs and outputs, draw the data
flow, calculate a toy case, identify its assumptions, predict a failure mode,
and design a test that could prove your implementation wrong.
\end{engineernote}

\section*{Why the learning pattern is repeated}

The repeated chapter compass, worked examples, failure cards, and retrieval
prompts are deliberate rather than decorative.
Constructive alignment connects observable learning outcomes to practice and
assessment\textsuperscript{\cite{biggsEnhancingTeachingConstructive1996}}.
Worked examples reduce avoidable search for novices, while support should fade
as expertise grows\textsuperscript{\cite{swellerCognitiveLoadTheory2019}}.
Practice testing and distributed practice have broad evidence for improving
retention\textsuperscript{\cite{dunloskyImprovingStudentsLearning2013}}, and
retrieval can outperform additional study after a delay
\textsuperscript{\cite{roedigerTestEnhancedLearning2006}}.

Words and explanatory graphics are placed together when the visual carries a
real part of the explanation.
This follows multimedia, signalling, coherence, and spatial-contiguity
principles; irrelevant illustration is omitted because visual density can add
extraneous processing instead of understanding
\textsuperscript{\cite{mayerUsingMultimediaElearning2017}}.

\begin{table}[htbp]
\centering
\caption{Evidence-informed learning patterns used throughout the book.}
\label{tab:evidence-informed-learning-patterns}
\begin{tblr}{
    width=\textwidth,
    colspec={Q[l,0.2\textwidth]X[l]X[l]},
    row{1}={font=\bfseries,bg=\maincolor!15!white},
    hlines,
    vlines}
Pattern & Purpose & How to use it \\
Chapter compass & Activate prerequisites and expose relationships & Predict
the four-stage chain before reading; redraw it afterwards \\
Worked example & Show the reasoning path, not only the answer & Explain why
each step is valid, then complete a partially worked variant \\
Failure card & Build diagnostic knowledge & Predict the symptom, locate the
violated assumption, and propose a discriminating test \\
Retrieval prompt & Strengthen later recall and expose gaps & Answer without
looking, check with feedback, and revisit after a delay \\
Transfer task & Prevent knowledge from remaining example-bound & Change the
domain, constraints, or data-generating process and adapt the method \\
Adjacent visual & Reduce split attention for spatial or process ideas & Trace
the diagram while explaining the corresponding text in your own words \\
\end{tblr}
\end{table}

These effects depend on prior knowledge, task complexity, feedback, and study
delay.
The book therefore offers a route rather than claiming that one layout works
identically for every reader.
Experienced readers may skip elementary scaffolds, but should still attempt
the failure and transfer prompts.

\section*{Evidence and verification policy}

The bibliography is not decoration.
Each citation should be close to the claim it supports, and important claims
should be triangulated where practical.

\begin{description}
    \item[Level A --- primary or normative evidence:]
        peer-reviewed papers, original technical reports, standards, laws,
        and official datasets or model cards.
    \item[Level B --- authoritative synthesis:]
        established textbooks, systematic reviews, and recognised handbooks.
    \item[Level C --- operational evidence:]
        official software documentation and reproducible implementation notes.
    \item[Level D --- illustration:]
        tutorials, analogies, and examples that clarify an idea but do not
        establish a general scientific claim.
\end{description}

For equations, verification means more than finding the same symbols in two
books.
Dimensions, units, boundary cases, sign conventions, and a numerical toy
example must agree.
For empirical claims, the dataset, split, metric, uncertainty, hardware, and
software configuration belong to the claim.

\begin{cautionbox}[The AI Almanach is a map, not the territory]
AI changes quickly and application domains contain specialised knowledge that
no single volume can exhaust.
Use the AI Almanach to build durable mental models and verification habits, then
read the cited primary sources for decisions that matter.
\end{cautionbox}

\section*{Notation quick reference}

\begin{table}[htbp]
\centering
\caption{Core notation used across chapters.}
\label{tab:core-notation}
\begin{tblr}{
    width=\textwidth,
    colspec={Q[c,0.13\textwidth] X[l] X[l]},
    row{1}={font=\bfseries, bg=\maincolor!15!white},
    hlines,
    vlines}
Symbol & Meaning & Typical shape or range \\
$x$ or $\mathbf{x}$ & scalar or feature vector &
$\mathbf{x}\in\mathbb{R}^{d}$ \\
$X$ & design matrix or token representations & $X\in\mathbb{R}^{n\times d}$ \\
$y$ & observed target & scalar, class, sequence, or structured object \\
$\hat y$ & model prediction & same semantic space as the target \\
$\theta$ & trainable model parameters & vector or collection of tensors \\
$\mathcal{L}$ & loss or empirical objective & usually $\mathbb{R}$ \\
$p(y\mid x)$ & conditional probability & $[0,1]$, sums or integrates to one \\
$\nabla_{\theta}\mathcal{L}$ & gradient of loss with respect to parameters &
same shape as $\theta$ \\
$\mathbb{E}[X]$ & expectation of a random variable & units of $X$ \\
\end{tblr}
\end{table}

\begin{checkpoint}
Choose a model you already know.
Can you name its input, target, parameters, objective, evaluation unit, one
assumption, and one failure mode?
Any blank is a useful reading goal rather than a reason for panic.
\end{checkpoint}
